\documentclass[11pt,a4paper]{article}

\usepackage[a4paper,margin=18mm,headheight=16pt]{geometry}
\usepackage{fontspec}
\setmainfont{Arial}
\setmonofont{Menlo}[Scale=0.82]

\usepackage{xcolor}
\definecolor{Ink}{HTML}{172033}
\definecolor{Muted}{HTML}{596579}
\definecolor{Blue}{HTML}{2855D9}
\definecolor{BlueLight}{HTML}{EEF3FF}
\definecolor{Amber}{HTML}{A86200}
\definecolor{AmberLight}{HTML}{FFF7E8}
\definecolor{Green}{HTML}{237A57}
\definecolor{Red}{HTML}{B42318}
\definecolor{Rule}{HTML}{D9DFEA}
\definecolor{Paper}{HTML}{F5F7FB}

\usepackage{microtype}
\usepackage{etoolbox}
\usepackage{needspace}
\usepackage{parskip}
\setlength{\parindent}{0pt}
\setlength{\parskip}{5pt plus 1pt}
\usepackage{enumitem}
\setlist[itemize]{leftmargin=5mm,itemsep=2pt,topsep=3pt}
\setlist[enumerate]{leftmargin=7mm,itemsep=3pt,topsep=3pt}
\providecommand{\tightlist}{\setlength{\itemsep}{2pt}\setlength{\parskip}{0pt}}

\usepackage{titlesec}
\titleformat{\part}[block]{\huge\bfseries\color{Blue}}{}{0pt}{}
\titlespacing*{\part}{0pt}{0pt}{14pt}
\titleformat{\section}{\Large\bfseries\color{Ink}}{\thesection}{0.7em}{}
\titleformat{\subsection}{\large\bfseries\color{Ink}}{\thesubsection}{0.65em}{}
\titlespacing*{\section}{0pt}{18pt}{7pt}
\titlespacing*{\subsection}{0pt}{13pt}{5pt}
\setcounter{secnumdepth}{2}
\setcounter{tocdepth}{2}

\usepackage{calc}
\usepackage{booktabs,tabularx,array,longtable}
\newcolumntype{Y}{>{\raggedright\arraybackslash}X}
\renewcommand{\arraystretch}{1.25}

\usepackage[most]{tcolorbox}
\tcbset{enhanced,boxrule=0pt,arc=1.5mm,left=4mm,right=4mm,top=3mm,bottom=3mm}
\newtcolorbox{keybox}[1][]{colback=BlueLight,colframe=BlueLight,
  borderline west={2pt}{0pt}{Blue},
  before upper={\ifstrempty{#1}{}{{\bfseries #1}\par\vspace{1pt}}}}
\newtcolorbox{warnbox}[1][]{colback=AmberLight,colframe=AmberLight,
  borderline west={2pt}{0pt}{Amber},
  before upper={\ifstrempty{#1}{}{{\bfseries #1}\par\vspace{1pt}}}}
\newtcolorbox{plainbox}[1][]{colback=Paper,colframe=Paper,
  borderline west={2pt}{0pt}{Rule},
  before upper={\ifstrempty{#1}{}{{\bfseries #1}\par\vspace{1pt}}}}

\usepackage{fontawesome5}
\newcommand{\yes}{\textcolor{Green}{\faCheck}}
\newcommand{\no}{\textcolor{Red}{\faTimes}}
\newcommand{\errdot}{\textcolor{Red}{\faCircle}}
\newcommand{\warndot}{\textcolor{Amber}{\faCircle}}

\usepackage{fancyhdr}
\pagestyle{fancy}
\fancyhf{}
\lhead{\small\color{Muted}OpenXR in Unity — Setup and Workflow Guide}
\rhead{\small\color{Muted}\thepage}
\renewcommand{\headrulewidth}{0.4pt}
\renewcommand{\headrule}{\hbox to\headwidth{\color{Rule}\leaders\hrule height \headrulewidth\hfill}}

\usepackage{xurl}
\usepackage[hidelinks]{hyperref}
\hypersetup{colorlinks=true,linkcolor=Blue,urlcolor=Blue}
\usepackage{bookmark}

\begin{document}

% ---------------------------------------------------------------------------
% Title page
% ---------------------------------------------------------------------------
\begin{titlepage}
\pagecolor{Ink}
\color{white}
\vspace*{42mm}
{\fontsize{14}{17}\selectfont\bfseries COURSE SETUP GUIDE\par}
\vspace{7mm}
{\fontsize{34}{39}\selectfont\bfseries OpenXR in Unity\par}
\vspace{5mm}
{\fontsize{17}{22}\selectfont Set up once, then two workflows for the rest of the
semester\par}
\vspace{16mm}
\begin{tcolorbox}[colback=white!8!Ink,colframe=white!18!Ink,boxrule=.6pt,arc=2mm]
\color{white}
\textbf{Meta Quest 2 / 3 / 3S. Windows and macOS.}\\[2mm]
This course is OpenXR only — Unity's own vendor-neutral packages. Not the Meta XR
All-in-One SDK, and not Building Blocks.
\end{tcolorbox}
\vfill
{\large Semester 2, 2026\par}
\vspace{4mm}
{\small\color{Rule}Guide version 1.0\par}
\end{titlepage}
\nopagecolor
\color{Ink}

\tableofcontents
\clearpage

This guide has two halves.

\textbf{Part 1 is setup.} You do it once, and then you leave it alone. It ends with a scene of
your own running standalone on a headset.

\textbf{Part 2 is workflows.} These are the two ways you will test your work, week after week,
for the rest of the course. You will come back to this half constantly.

Everything here assumes you have already worked through \texttt{Software\_and\_Frameworks.md} and
the Unity + GitHub guide, that you have Unity 6.3 LTS with Android Build Support
installed, and that you have a Unity project sitting inside a working repository. It
assumes no prior XR experience at all.

\begin{warnbox}[One warning before you start]
This course uses \textbf{OpenXR only} --- Unity's own
packages on the open standard. We do \textbf{not} use the Meta XR All-in-One SDK and we do
\textbf{not} use Building Blocks. Most Quest tutorials you find online do use them. Those
tutorials aren't wrong, they're just a different stack, and installing both at once
will break your project in ways that are genuinely hard to diagnose. If a tutorial
starts by importing the Meta XR SDK, close it.
\end{warnbox}


\clearpage
\part*{Part 0 --- Before you start}

\section*{Prerequisites}\label{prerequisites}
\addcontentsline{toc}{section}{Prerequisites}

You need all of these before section 1. Every one of them is covered in an earlier guide.

\begin{itemize}
\tightlist
\item
  Unity Hub installed, and \textbf{Unity 6.3 LTS} (\texttt{6000.3.x}) installed through it
\item
  Unity installed with \textbf{Android Build Support}, including \textbf{OpenJDK} and
  \textbf{Android SDK \& NDK Tools}
\item
  A course repository with the supplied \texttt{.gitignore} already in place
\item
  A Unity project inside that repository that opens without errors
\item
  A Meta Quest 2, 3, or 3S, charged, with its controllers paired
\item
  \textbf{A USB-C cable that carries data.} Charge-only cables look identical and will waste
  an afternoon. Everything in this guide is done tethered
\end{itemize}

\section*{The two workflows}\label{the-two-workflows}
\addcontentsline{toc}{section}{The two workflows}

There are two ways you will test your work, and both run identically on Windows and
macOS. You need both --- they answer different questions.

\begin{longtable}[]{@{}
  >{\raggedright\arraybackslash}p{(\linewidth - 6\tabcolsep) * \real{0.2500}}
  >{\raggedright\arraybackslash}p{(\linewidth - 6\tabcolsep) * \real{0.2500}}
  >{\raggedright\arraybackslash}p{(\linewidth - 6\tabcolsep) * \real{0.2500}}
  >{\raggedright\arraybackslash}p{(\linewidth - 6\tabcolsep) * \real{0.2500}}@{}}
\toprule\noalign{}
\begin{minipage}[b]{\linewidth}\raggedright
Workflow
\end{minipage} & \begin{minipage}[b]{\linewidth}\raggedright
What it's for
\end{minipage} & \begin{minipage}[b]{\linewidth}\raggedright
Speed
\end{minipage} & \begin{minipage}[b]{\linewidth}\raggedright
Headset
\end{minipage} \\
\midrule\noalign{}
\endhead
\bottomrule\noalign{}
\endlastfoot
\textbf{A --- Build to the headset} & The truth. Real performance, real tracking, real comfort. & Minutes & Required \\
\textbf{B --- XR Device Simulator} & Fast iteration in the Editor with keyboard and mouse. & Seconds & Not required \\
\end{longtable}

A third option, \textbf{Quest Link}, exists on Windows only and is described briefly in
section 8. We don't teach it, because these two cover everything and work on every
machine on the course.


\clearpage
\part*{Part 1 --- Setup}

\section{What OpenXR actually is}\label{what-openxr-actually-is}

For most of VR's short history, every headset came with its own SDK. Code written for a
Rift didn't run on a Vive. Vendors competed by locking developers in, and developers paid
for it every time they wanted to ship somewhere new.

\textbf{OpenXR} is the industry's answer to that. It's an open standard, managed by the Khronos
Group --- the same people behind Vulkan and OpenGL --- that defines a common language between
an application and a headset. Your app says ``give me the pose of the right hand'' or ``tell
me when the trigger is pressed'', and whatever runtime is underneath answers in the same
way, whether that's a Quest, a Vive, a Valve Index, or a headset that doesn't exist yet.

The stack looks like this:

\begin{verbatim}
    Your Unity scene
          ↓
    XR Interaction Toolkit     ← grabbing, teleporting, UI. Unity's, vendor-neutral.
          ↓
    Unity OpenXR Plugin        ← translates Unity's XR API into OpenXR calls
          ↓
    OpenXR runtime             ← supplied by Meta on Quest, by Valve on Index, etc.
          ↓
    The headset
\end{verbatim}

You write against the top two layers. The bottom two swap out underneath you.

\textbf{The Oculus XR Plugin.} Older tutorials use a Unity package called the Oculus XR
Plugin, which was the pre-OpenXR path for Quest development. It is deprecated. Don't
install it --- and treat the rest of any tutorial that recommends it with suspicion, since
it predates the current toolchain entirely.

\textbf{The Meta XR SDK and Building Blocks.} Meta's own SDK is current, well maintained, and
genuinely good; its Building Blocks drop a working camera rig or hand tracking into a
scene in seconds, and most Quest tutorials online use them. We don't, for two reasons.
What you learn on OpenXR and the XR Interaction Toolkit transfers to any headset, while
Building Block prefab names transfer nowhere. And when a Building Block works you learn
nothing from it, whereas the rig you assemble by hand in section 6 is one you can debug
when it breaks.

\begin{warnbox}[Do not install both stacks]
The Meta SDK and Unity's OpenXR packages both want to
configure the same project settings, own the same camera, and drive the same input.
Installing the Meta SDK ``just to try a tutorial'' and then removing it leaves settings
behind. Symptoms are listed in \hyperref[troubleshooting]{Section~13}.
\end{warnbox}


\section{Headset and computer preparation}\label{headset-and-computer-preparation}

\textbf{The course headsets are already set up.} Developer Mode is enabled on all of them, and
they are ready to receive your builds. You do not need a Meta developer account, an
organisation, or the Meta Horizon mobile app to work on this course.

There is one thing you do have to do, and it is per-computer:

\textbf{Authorise your computer.} Plug the headset into your machine with the USB-C data cable
and put the headset on. A prompt appears \emph{inside the headset} asking whether to allow USB
debugging. Tick \textbf{Always allow from this computer} and accept.

You will not see this prompt on your monitor. It is in the headset. Students routinely sit
staring at a failed build for ten minutes with an unanswered dialogue floating in front of
a headset on the desk beside them.

You will need to do this again on each new computer you work at, and sometimes again after
a headset system update.

\begin{plainbox}[Using your own headset]
If you're working on a personal Quest rather than a course
one, you'll need to enable Developer Mode on it yourself first: create a Meta developer
account, create an organisation on the developer dashboard, then turn on Developer Mode
for your headset in the Meta Horizon mobile app and restart the headset. Links are in
\texttt{Software\_and\_Frameworks.md}.
\end{plainbox}

\begin{keybox}[Checkpoint]
Open Meta Quest Developer Hub (MQDH) and check that your headset appears
as connected over USB. If it does, you're ready. If it doesn't, fix it now --- nothing
later in this guide will work until this does.
\end{keybox}


\section{Switching the project to Android}\label{switching-the-project-to-android}

The Quest runs Android. Before Unity can build anything for it, the project has to be
switched over.

\subsection{Switch platform}\label{switch-platform}

\textbf{File → Build Profiles} → in the \textbf{Platforms} list, select \textbf{Android} →
\textbf{Switch Platform}.

This re-imports every asset in the project into Android-compatible formats. On an empty
project it takes a moment. On a project full of textures it can take a long while. Do it
early, while the project is still small.

Set \textbf{Texture Compression} to \textbf{ASTC}, which is what Quest hardware wants.

\subsection{Player settings that matter}\label{player-settings-that-matter}

\textbf{Edit → Project Settings → Player}, with the Android tab selected.

Under \textbf{Other Settings}:

\begin{longtable}[]{@{}
  >{\raggedright\arraybackslash}p{(\linewidth - 4\tabcolsep) * \real{0.3333}}
  >{\raggedright\arraybackslash}p{(\linewidth - 4\tabcolsep) * \real{0.3333}}
  >{\raggedright\arraybackslash}p{(\linewidth - 4\tabcolsep) * \real{0.3333}}@{}}
\toprule\noalign{}
\begin{minipage}[b]{\linewidth}\raggedright
Setting
\end{minipage} & \begin{minipage}[b]{\linewidth}\raggedright
Value
\end{minipage} & \begin{minipage}[b]{\linewidth}\raggedright
Why
\end{minipage} \\
\midrule\noalign{}
\endhead
\bottomrule\noalign{}
\endlastfoot
\textbf{Company Name} / \textbf{Product Name} & Anything sensible & These generate your package name --- set them before your first build \\
\textbf{Graphics APIs} & \textbf{Vulkan} & The modern path on Quest. Remove OpenGLES3 if it's listed above Vulkan \\
\textbf{Scripting Backend} & \textbf{IL2CPP} & Required --- Mono won't build for Quest \\
\textbf{Target Architectures} & \textbf{ARM64} only & Untick ARMv7. Quest hardware is 64-bit and builds are smaller without it \\
\textbf{Minimum API Level} & \textbf{Android 12L (API 32)} & Meta's floor for current Horizon OS \\
\textbf{Active Input Handling} & \textbf{Input System Package (New)} & The Interaction Toolkit needs it. Changing this restarts the Editor \\
\end{longtable}

Under \textbf{Resolution and Presentation}, tick \textbf{Optimized Frame Pacing}.

\subsection{Quality settings}\label{quality-settings}

\textbf{Edit → Project Settings → Quality}. Standalone VR is a mobile GPU rendering two eyes
at high refresh rate, and it will punish you for treating it like a desktop.

Pick a low or medium quality level as the Android default, and inside it keep shadows
modest and anti-aliasing at \textbf{4x MSAA} (MSAA is comparatively cheap on this hardware and
matters a lot for how solid edges look up close).

Don't over-tune this now. Get a build running first, then optimise against real numbers
from a real device.


\section{Installing the XR packages}\label{installing-the-xr-packages}

All of these come from Unity's own package registry, inside the project. Nothing here is
downloaded from a website or imported from the Asset Store.

Install in this order --- each one expects the previous.

\subsection{XR Plugin Management}\label{xr-plugin-management}

\textbf{Edit → Project Settings → XR Plug-in Management} → \textbf{Install XR Plugin Management}.

This is the settings panel where you choose an XR backend per platform. It doesn't do
anything by itself.

\subsection{Unity OpenXR Plugin}\label{unity-openxr-plugin}

Still in \textbf{XR Plug-in Management}, select the \textbf{Android} tab (the little robot icon) and
tick \textbf{OpenXR}. Unity installs the OpenXR plugin for you.

You'll see a red warning icon appear next to it. That's expected --- it's Project Validation
telling you OpenXR isn't configured yet. Section 5 fixes it.

\subsection{XR Interaction Toolkit}\label{xr-interaction-toolkit}

\textbf{Window → Package Manager} → \textbf{Unity Registry} → search \textbf{XR Interaction Toolkit} →
\textbf{Install}.

Then, still in Package Manager with the toolkit selected, open the \textbf{Samples} tab and
import:

\begin{itemize}
\tightlist
\item
  \textbf{Starter Assets} --- this contains the default input action maps that bind physical
  controller buttons to named actions. You need it. Without it you have a camera rig with
  no input.
\item
  \textbf{XR Device Simulator} --- the fast-iteration workflow from section 11. Import it now
  while you're here.
\end{itemize}

\subsection{Later additions}\label{later-additions}

You don't need these yet. They arrive with the activities that use them, and the
activity brief will tell you.

\begin{itemize}
\tightlist
\item
  \textbf{XR Hands} (\texttt{com.unity.xr.hands}) --- hand tracking instead of controllers
\item
  \textbf{AR Foundation} + \textbf{Unity OpenXR: Meta} --- passthrough and mixed reality on Quest 3
  and 3S
\end{itemize}

\begin{keybox}[Commit now]
Your \texttt{Packages/manifest.json} and \texttt{Packages/packages-lock.json} have
changed. Those two files are how a different machine reconstructs exactly this package
set, so they belong in Git. Commit before you go further.
\end{keybox}


\section{Configuring OpenXR}\label{configuring-openxr}

Installing OpenXR isn't the same as configuring it. This section is where most setup
failures actually happen, and it's worth going slowly.

\subsection{Enable the Meta feature group}\label{enable-the-meta-feature-group}

\textbf{Edit → Project Settings → XR Plug-in Management} → \textbf{Android} tab.

Under the OpenXR entry you'll find a list of \textbf{OpenXR Feature Groups}. Tick \textbf{Meta Quest
Support}.

This is what tells the OpenXR plugin to produce a build the Quest runtime will accept ---
the correct Android manifest entries, the right permissions, the device declaration.
Without it your app will build happily and then refuse to launch.

\subsection{Interaction profiles}\label{interaction-profiles}

\textbf{Edit → Project Settings → XR Plug-in Management → OpenXR} → \textbf{Android} tab.

Find \textbf{Enabled Interaction Profiles} and add:

\begin{itemize}
\tightlist
\item
  \textbf{Oculus Touch Controller Profile}
\end{itemize}

An interaction profile is OpenXR's description of a specific physical controller --- which
buttons exist, where they are, what they're called. Your input actions bind to profiles.
\textbf{If this list is empty, your controllers will track as objects in space but no button
will do anything}, which is a maddening bug to debug from scratch.

If you're on Quest 3 or 3S you can add \textbf{Meta Quest Touch Plus Controller Profile} as
well. Adding profiles you don't own is harmless.

\subsection{Render mode}\label{render-mode}

On the same page, set \textbf{Render Mode} to \textbf{Single Pass Instanced}.

VR renders the scene twice, once per eye. Single Pass Instanced does that in one pass
instead of two and is close to free performance. Leave it on unless something specific
breaks.

\subsection{Project Validation}\label{project-validation}

\textbf{Edit → Project Settings → XR Plug-in Management → Project Validation}.

This panel checks your project against everything OpenXR and the Meta feature group
expect, and it is the single most useful diagnostic tool in this whole guide. Get into the
habit of coming back here whenever something is strange.

Each issue has a \textbf{Fix} button. Most can be fixed automatically. Work top to bottom until
the Android tab is clean.

A note on the two icons:

\begin{itemize}
\tightlist
\item
  \errdot~\textbf{Errors} --- must be fixed. Your build will fail or misbehave.
\item
  \warndot~\textbf{Warnings} --- should usually be fixed, but read them. Some are recommendations
  about quality settings that you may have deliberately chosen otherwise.
\end{itemize}

\begin{keybox}[Checkpoint]
Android tab of Project Validation shows no red errors. If it does, don't
continue --- nothing in section 6 will work, and you'll waste time debugging a scene when
the problem is in settings.
\end{keybox}


\section{The known-good scene}\label{the-known-good-scene}

Now you build the smallest possible scene that proves the whole stack works. It has a
floor, a thing to look at, and a camera rig you can move your head inside.

It is deliberately trivial. That's the point --- \textbf{when this scene fails, the failure is in
your setup and nothing else}. Keep it in your project all semester. Every time something
mysterious breaks in a real activity, run this scene first to find out whether the problem
is your work or your machine.

You will also run this exact scene through both workflows in Part 2, so that you know what
``working'' looks like in each of them before you have real work on the line.

\subsection{Create the scene}\label{create-the-scene}

\textbf{File → New Scene} → choose the \textbf{Basic (URP)} template → save it as
\texttt{KnownGood.unity} somewhere sensible, like \texttt{Assets/Scenes/}.

Delete the \textbf{Main Camera} that came with the template. The XR rig brings its own.

\subsection{Build the XR rig}\label{build-the-xr-rig}

\textbf{GameObject → XR → XR Origin (VR)}.

Look at what that just created in the Hierarchy, because this is the part worth
understanding:

\begin{verbatim}
XR Origin (XR Rig)          ← the player's position in the world. Move THIS to move the player.
└── Camera Offset           ← height compensation, managed for you
    ├── Main Camera         ← the headset. Its transform is driven by tracking — never set it yourself.
    ├── Left Controller     ← driven by the left controller's tracked pose
    └── Right Controller
\end{verbatim}

The mental model that saves the most confusion later: \textbf{the camera's position is an
output, not an input.} It reports where the player's head is. If you want to move the
player, you move the XR Origin, and the head keeps its offset within it. Trying to set the
camera's transform directly is the most common beginner mistake in XR, and it fails in a
way that makes people motion sick.

An \textbf{XR Interaction Manager} object also appeared. It's the broker between things that
can interact and things that can be interacted with. One per scene, and you can ignore it.

Select \textbf{XR Origin} and set \textbf{Tracking Origin Mode} to \textbf{Floor}. This means y = 0 is
the player's real floor, and a 1.7-metre-tall object will look 1.7 metres tall. On
\textbf{Device} mode, y = 0 is wherever their head happened to be, and your whole scene will
feel like it's floating.

\subsection{Wire up input}\label{wire-up-input}

Create an empty GameObject, name it \texttt{Input\ Action\ Manager}, and add the \textbf{Input Action
Manager} component to it.

In its \textbf{Action Assets} list, add one element and assign \textbf{XRI Default Input Actions} ---
it came from the Starter Assets sample you imported in 4.3, under
\texttt{Assets/Samples/XR\ Interaction\ Toolkit/.../Starter\ Assets/}.

This component's only job is to enable those input actions at runtime. It is one component
doing one small thing, and forgetting it produces a rig where everything tracks correctly
and no button works --- the same symptom as a missing interaction profile, which is why
these two are worth remembering as a pair.

\subsection{Something to stand on and something to look at}\label{something-to-stand-on-and-something-to-look-at}

\begin{itemize}
\tightlist
\item
  \textbf{GameObject → 3D Object → Plane}, position \texttt{(0,\ 0,\ 0)}. This is your floor.
\item
  \textbf{GameObject → 3D Object → Cube}, position \texttt{(0,\ 1.2,\ 2)}, scale \texttt{(0.3,\ 0.3,\ 0.3)}.
\end{itemize}

The cube is at roughly chest height, two metres in front of you. Its job is to give you
something with an obvious real-world size, so that when you put the headset on you can
immediately tell whether the scale of the world is right.

Give both a material with some colour if you like. Flat grey is harder to judge depth
against than you'd expect.

\subsection{Add the scene to the build}\label{add-the-scene-to-the-build}

\textbf{File → Build Profiles} → \textbf{Scene List} → \textbf{Add Open Scenes}, and make sure
\texttt{KnownGood} is ticked and at the top.

A scene that isn't in this list will not be in your build. ``Black screen on launch'' is
very often just this.

Save the scene. Commit.


\section{Next time: start from the VR template}\label{next-time-start-from-the-vr-template}

Everything in sections 3 to 6 was you adding XR to a project that didn't have it. That's
the version worth doing once, because it's the version where you can see what each setting
is for.

For your next project, you don't have to do any of it. In \textbf{Unity Hub → New Project},
choose the \textbf{VR} template. It arrives with XR Plugin Management installed, OpenXR enabled
for Android, the XR Interaction Toolkit and its Starter Assets already imported, and a
sample scene with a working rig --- the end state of Part 1, ready to go.

Two things to know about it:

\textbf{Check the settings anyway.} The template configures itself for a general OpenXR target,
not specifically for Quest. Walk section 5 again --- the \textbf{Meta Quest Support} feature
group and the \textbf{Oculus Touch Controller Profile} in particular --- and run \textbf{Project
Validation} before you build. It takes two minutes and it's the difference between a
build that runs and a black screen.

\textbf{You still need Android.} The template doesn't know you're targeting a headset that runs
Android, so section 3 still applies: switch platform, then the player settings.

Use the template from here on. Sections 3 to 6 are now your reference for when something
in it looks wrong, and for understanding what the template did on your behalf.

\textbf{Part 1 is done.} You have a configured project and a scene worth testing. Everything
from here is about getting it in front of your eyes.


\clearpage
\part*{Part 2 --- Workflows}

\section{Choosing a workflow}\label{choosing-a-workflow}

You now have two ways to see your work. They are not alternatives to each other --- they
are tools for different moments, and a good week uses both.

\begin{longtable}[]{@{}lll@{}}
\toprule\noalign{}
& \textbf{A --- On the headset} & \textbf{B --- XR Device Simulator} \\
\midrule\noalign{}
\endhead
\bottomrule\noalign{}
\endlastfoot
\textbf{Speed per iteration} & Minutes & Seconds \\
\textbf{Headset required} & Yes & No \\
\textbf{Real performance} & Yes & No \\
\textbf{Real comfort and scale} & Yes & No \\
\textbf{Real tracking and controllers} & Yes & No \\
\textbf{Windows and macOS} & Yes & Yes \\
\end{longtable}

\subsection{The weekly loop}\label{the-weekly-loop}

\begin{enumerate}
\def\labelenumi{\arabic{enumi}.}
\tightlist
\item
  \textbf{Build the thing in the simulator.} Logic, layout, interactions, does the button do
  what it should. This is where most of your hours go, and it needs no headset --- which
  means you can work anywhere, without booking equipment.
\item
  \textbf{Confirm on the headset before you call it done.} Every session, not just at the end
  of a project. Scale, comfort, reachability, and framerate are all things the simulator
  cannot tell you, and all things that can invalidate an afternoon's work.
\end{enumerate}

\textbf{Work that has only ever run in the simulator isn't finished work.} If you take one
thing from Part 2, take that.

\subsection{A note on Quest Link}\label{a-note-on-quest-link}

On Windows, Meta's desktop software offers \textbf{Quest Link}, which runs your scene live in
the headset when you press Play --- no build step, real tracking. If you have a Windows
machine and want to try it, install the Meta Horizon desktop app, connect by cable, set
Meta Quest Link as the active OpenXR runtime, and enable OpenXR on the \textbf{Windows, Mac,
Linux} tab of XR Plug-in Management as well as the Android tab (Play mode reads the
desktop settings, not the Android ones --- that's the step people miss).

We don't teach it, for two reasons: it doesn't exist on macOS, and its performance is a
flattering lie, because your PC's graphics card is doing the rendering rather than the
headset's chip. It can replace step 1 of the loop above. It cannot replace step 2.


\section{Workflow A --- Building and running on the headset}\label{workflow-a-building-and-running-on-the-headset}

This is the ground truth. Everything else is an approximation of it.

\subsection{Your first build}\label{your-first-build}

\begin{enumerate}
\def\labelenumi{\arabic{enumi}.}
\tightlist
\item
  Connect the headset by USB-C and put it on briefly to confirm there's no pending
  authorisation prompt (section 2).
\item
  \textbf{File → Build Profiles} → \textbf{Android} → check your device is listed under
  \textbf{Run Device}. Hit the refresh button beside it if it isn't.
\item
  Click \textbf{Build and Run}.
\item
  Choose an output folder. \textbf{Put it outside your repository} --- a folder called \texttt{Builds}
  next to the project, not inside it. Builds are large, they change completely every
  time, and they must never be committed. The course \texttt{.gitignore} covers the usual
  locations, but the reliable habit is keeping them out of the repo entirely.
\item
  Wait.
\end{enumerate}

\subsection{What ``wait'' means}\label{what-wait-means}

\textbf{Your first build will take a long time.} Ten to twenty minutes is normal, and on an
older laptop it can be more. IL2CPP is converting your C\# to C++ and then compiling it for
ARM64, and there's no shortcut.

Later builds of the same project are much faster --- typically one to three minutes ---
because most of that work is cached.

Plan around this. Get one build through early in the semester, so that the twenty-minute
version is behind you rather than in front of you on a day you need to show someone
something.

\subsection{Finding your app on the headset}\label{finding-your-app-on-the-headset}

\textbf{Build and Run} launches the app for you. But once you take the headset off and go back
to the home screen, your app is not with the ones you bought.

In the headset: open the \textbf{Library}, then find the source or category filter, and choose
\textbf{Unknown Sources}. Everything you have side-loaded lives there, listed by the Product
Name you set back in section 3.2 --- which is why setting it to something recognisable
mattered.

\subsection{The repeat loop}\label{the-repeat-loop}

Once the first build is through, your normal cycle is: change something in Unity →
\textbf{Build and Run} → put the headset on. If your device is already connected and
authorised, that's a couple of minutes.

If the headset is connected and Unity can't see it, the fix is almost always in
\hyperref[troubleshooting]{Section~13} under \emph{headset not detected}.

\begin{keybox}[Checkpoint]
Your cube is floating in front of you at roughly chest height, it stays
put when you move your head, and both controllers appear where your hands are. If so,
your entire setup is correct. This is the moment Part 1 was for.
\end{keybox}


\section{Meta Quest Developer Hub}\label{meta-quest-developer-hub}

MQDH is Meta's desktop companion for the headset. It isn't a third workflow --- it sits
alongside both of them, and it runs on \textbf{Windows and macOS} alike.

You'll see it called \textbf{ODH} or \textbf{Oculus Developer Hub} in older documentation, and in
Meta's own URLs. Same tool, older name.

Connect over USB. MQDH also offers ADB over Wi-Fi, and on campus it is not worth the
trouble --- the wireless here is unreliable enough that a dropped connection mid-deploy will
cost you more time than the cable ever will.

\subsection{What you'll use it for}\label{what-youll-use-it-for}

\textbf{Deploying builds.} Drag an \texttt{.apk} onto the \textbf{Device Manager} and it installs. Useful
when Unity's Build and Run is being uncooperative, and useful for handing a build to a
teammate who doesn't have your project.

\textbf{Mirroring the headset to your screen.} Start casting and what the wearer sees appears
in a window on your computer:

\begin{itemize}
\tightlist
\item
  Your tutor can see what you're seeing without wrestling the headset off your head
\item
  Your team can watch a playtest and take notes at the same time
\item
  You can spot in a second that the problem is a controller not tracking
\end{itemize}

\textbf{Capturing video and screenshots.} Recording straight from the device beats filming a
screen with your phone, whenever you need to show what your work actually looks like.

\textbf{Reading logs.} When a build launches and immediately dies, the reason is in the log,
and this is where you read it.

\subsection{Set it up now, not later}\label{set-it-up-now-not-later}

Install MQDH, connect your headset, and confirm casting works --- before you need it under
pressure in a lab. It takes five minutes on a quiet afternoon.


\section{Workflow B --- The XR Device Simulator}\label{workflow-b-the-xr-device-simulator}

The simulator gives you a virtual headset and two virtual controllers driven by your
keyboard and mouse, running in the Editor's Play mode. No headset, no build, no cable.
It works identically on Windows and macOS.

This is where most of your development time should go.

\subsection{Setting it up}\label{setting-it-up}

You imported the sample in section 4.3. Find it under
\texttt{Assets/Samples/XR\ Interaction\ Toolkit/.../XR\ Device\ Simulator/}, and drag the simulator
prefab into your scene.

That's the whole setup. Press \textbf{Play}.

\subsection{Driving it}\label{driving-it}

An on-screen panel appears showing the current controls and which device you're currently
driving. Read it --- it's the authoritative reference, and it updates as you change modes.

The broad shape of it:

\begin{itemize}
\tightlist
\item
  \textbf{Mouse} moves the head, so you look around
\item
  \textbf{WASD} moves the rig through the world
\item
  Modifier keys hand control over to the \textbf{left} or \textbf{right} controller, so mouse and
  keyboard drive that hand instead of your head
\item
  While driving a controller, other keys press its trigger, grip, and buttons
\end{itemize}

Spend two minutes with the on-screen panel the first time. Once the ``which thing am I
currently driving'' idea clicks, the rest is muscle memory.

\subsection{What the simulator will not tell you}\label{what-the-simulator-will-not-tell-you}

This section matters more than the rest of this workflow. The simulator is a model of a
headset, and every model leaves things out. It is silent --- not wrong, \emph{silent} --- about:

\textbf{Performance.} Your Editor is running on a desktop GPU. The Quest runs on a mobile chip
and needs to hold 72 frames a second in stereo. A scene that runs beautifully in the
simulator can be unusable on device, and the simulator will never warn you.

\textbf{Comfort.} Motion sickness is the single most common failure in student XR work, and
it's invisible on a monitor. Snap turns, acceleration, moving the horizon, poorly placed
UI --- you cannot evaluate any of it here.

\textbf{Scale.} Objects look right on screen and wrong at real size, constantly. A doorway that
looks fine in the Scene view can be intimidatingly small when you're standing in it.

\textbf{Real tracking.} Real hands move differently to a mouse. Real controllers occlude,
drift, and leave tracking volume. Real people can't reach behind themselves.

\textbf{Reachability and ergonomics.} Whether the player can comfortably grab the thing you
want them to grab, in the position a real body is actually in.

\begin{keybox}
Use the simulator to find out \textbf{whether it works}. Use the headset to find out
\textbf{whether it's good}. Those are different questions.
\end{keybox}


\section{Working across two machines}\label{working-across-two-machines}

Most of you will work on a lab PC some days and your own machine on others. The Unity +
GitHub guide covers moving a project between computers; this is the XR-specific part of
that.

\subsection{What has to match}\label{what-has-to-match}

\begin{longtable}[]{@{}
  >{\raggedright\arraybackslash}p{(\linewidth - 2\tabcolsep) * \real{0.5000}}
  >{\raggedright\arraybackslash}p{(\linewidth - 2\tabcolsep) * \real{0.5000}}@{}}
\toprule\noalign{}
\endhead
\bottomrule\noalign{}
\endlastfoot
\textbf{Unity version} & Exactly \texttt{6000.3.x}, the same as everyone else. Opening a project in a newer Editor upgrades it irreversibly for whoever opens it next \\
\textbf{Package versions} & Handled for you, as long as \texttt{Packages/manifest.json} and \texttt{Packages/packages-lock.json} are committed \\
\textbf{Project settings} & Committed, as long as \texttt{ProjectSettings/} is in the repo --- it is, in the supplied \texttt{.gitignore} \\
\end{longtable}

\subsection{What is per-machine and doesn't travel}\label{what-is-per-machine-and-doesnt-travel}

None of this is in Git, and none of it should be. Expect to redo it the first time you sit
at a new machine:

\begin{itemize}
\tightlist
\item
  Signing in to Unity Hub and applying your licence
\item
  Installing and signing in to MQDH
\item
  Authorising USB debugging, which is per computer \emph{and} per headset (section 2)
\end{itemize}

\subsection{What must never be committed}\label{what-must-never-be-committed}

\begin{itemize}
\tightlist
\item
  \texttt{Library/}, \texttt{Temp/}, \texttt{Logs/}, \texttt{obj/} --- Unity regenerates all of it
\item
  \textbf{Builds and \texttt{.apk} files} --- large, binary, and completely disposable. Keep your build
  output folder outside the repository entirely
\item
  Anything under \texttt{Assets/Samples/} that you haven't modified is fine to commit, and
  simplest to leave in --- it's small, and it keeps everyone's input actions identical
\end{itemize}

\subsection{The habit that protects you}\label{the-habit-that-protects-you}

\textbf{Make a device build at the end of every session}, whichever machine you're on. It costs
two minutes, it confirms the project still builds from a clean-ish state, and it means you
find out that Wednesday's work doesn't hold framerate on Wednesday.


\clearpage
\part*{Part 3 --- When it goes wrong}

\section{Troubleshooting}\label{troubleshooting}

One symptom per heading. Find yours.

\subsection{Setup and build}\label{setup-and-build}

\textbf{Unity says ``no Android SDK found'' or ``Android SDK not found at path''}
The Android modules aren't installed, or Unity has lost track of them. Unity Hub →
\textbf{Installs} → three-dot menu on 6.3 → \textbf{Add modules} → confirm \textbf{Android Build
Support}, \textbf{OpenJDK}, and \textbf{Android SDK \& NDK Tools} are all ticked. If they are,
check \textbf{Edit → Preferences → External Tools} and re-tick the ``use recommended version''
boxes for SDK, NDK, and JDK.

\textbf{The build fails with errors mentioning IL2CPP or ARMv7}
Check section 3.2: \textbf{Scripting Backend} must be IL2CPP and \textbf{Target Architectures} must
be ARM64 with ARMv7 unticked.

\textbf{Project Validation shows red errors I don't understand}
Press \textbf{Fix}. That's what it's there for. If a fix doesn't stick, note the exact wording
and bring it to your tutor --- a persistent validation error usually means two packages are
disagreeing, which is worth a second pair of eyes.

\subsection{Device connection}\label{device-connection}

\textbf{Unity or MQDH can't see the headset}
Work through these in order:
1. Is the cable a \textbf{data} cable? Charge-only USB-C cables are extremely common and look
identical. Try another.
2. Put the headset on and look for the USB debugging prompt (section 2).
3. Is the headset awake? A sleeping headset drops off the connection.
4. Try a different USB port --- prefer one directly on the machine over a hub or dock.
5. Restart the headset.

\textbf{It was working, and now Unity's Run Device list is empty}
Usually the authorisation was revoked by a headset update. Unplug, replug, put the headset
on, accept the prompt again.

\subsection{On the headset}\label{on-the-headset}

\textbf{The app builds successfully but I can't find it on the headset}
It's under \textbf{Library → Unknown Sources} (section 9.3), listed by Product Name, not by
project name.

\textbf{Black screen when the app launches}
In rough order of likelihood:
1. Your scene isn't in the \textbf{Scene List} in Build Profiles (section 6.5)
2. \textbf{Meta Quest Support} isn't ticked in the OpenXR feature groups (section 5.1)
3. OpenXR isn't ticked at all on the \textbf{Android} tab of XR Plug-in Management
4. Something threw an exception at startup --- read the log in MQDH

\textbf{The app runs but the controllers don't do anything}
Two causes, and they produce an identical symptom, so check both:
1. \textbf{Enabled Interaction Profiles} is empty (section 5.2)
2. The \textbf{Input Action Manager} is missing, or has no action asset assigned (section 6.3)

\textbf{The controllers don't appear at all, or float in the wrong place}
Confirm they're paired and charged. If they track in the headset's home environment but
not in your app, it's an interaction profile problem --- see above.

\textbf{Everything is the wrong size, or I'm standing underneath the floor}
\textbf{Tracking Origin Mode} is set to \textbf{Device} instead of \textbf{Floor} on the XR Origin
(section 6.2).

\textbf{The scene moves when I move my head, and it makes me feel ill}
Something is writing to the Main Camera's transform. Nothing in your code should ever do
that --- move the XR Origin instead (section 6.2).

\subsection{Simulator}\label{simulator}

\textbf{The simulator prefab is in the scene but nothing responds}
Check that \textbf{Active Input Handling} is set to \textbf{Input System Package (New)} (section
3.2). Also confirm the Game view has focus --- the simulator only receives input when the
Game view is the focused window. Click into it once.

\textbf{I can look around but the controllers never move}
You haven't taken control of a hand. Check the on-screen control panel for the modifier
keys that switch which device you're driving (section 11.2).

\subsection{The one that's hardest to spot}\label{the-one-thats-hardest-to-spot}

\textbf{Strange, contradictory errors after following an online tutorial}
If you installed the Meta XR All-in-One SDK at any point --- even briefly, even if you
removed it --- you may have two XR stacks fighting. Signs: duplicate or unfamiliar entries
in XR Plug-in Management, an \texttt{OVRManager} or \texttt{OVRCameraRig} anywhere in the project,
Project Validation errors that reappear after you fix them, or compile errors mentioning
\texttt{OVR} or \texttt{Oculus}.

The reliable fix is to remove the Meta SDK completely, delete the \texttt{Library} folder, and
reopen the project --- and if it's still confused, to check out a clean copy of the project
from Git and reapply your work. This is a strong argument for committing often.


\section{Where to get help}\label{where-to-get-help}

\textbf{Start here, in this order:}

\begin{enumerate}
\def\labelenumi{\arabic{enumi}.}
\tightlist
\item
  \textbf{Project Validation} --- \texttt{XR\ Plug-in\ Management\ →\ Project\ Validation}. Genuinely fixes
  a large share of problems by itself
\item
  \textbf{Section 13} above
\item
  \textbf{\texttt{Software\_and\_Frameworks.md}} --- every version number, package name, and official
  documentation link for the whole stack lives there, and it's the authority if anything
  in this guide contradicts it
\item
  \textbf{The known-good scene} from section 6 --- run it. It tells you instantly whether the
  problem is your work or your setup, and that halves the search space
\item
  \textbf{Your tutor} --- bring the exact error text and what you've already tried
\end{enumerate}

\textbf{Official documentation} is linked from \texttt{Software\_and\_Frameworks.md} rather than
repeated here, so that there's one place to keep it current: Unity's OpenXR Plugin and XR
Interaction Toolkit manuals, and Meta's device and MQDH documentation.

\textbf{A note on searching for help.} Most Quest tutorials and most Stack Overflow answers
assume the Meta XR SDK. When you search, add \texttt{OpenXR} and \texttt{XR\ Interaction\ Toolkit} to your
query, and be suspicious of any answer that tells you to install something from the Asset
Store. If an answer's first step is ``add a Building Block'', it's solving a different
problem than the one you have.

\textbf{If a setup problem isn't answered anywhere in these three guides, that's a gap in the
guides --- tell your tutor so it can be fixed for everyone.}

\end{document}
